\begin{frame}{1. 研究背景：检索器的服务对象从人转向 Agent}
    \begin{center}
        \keynote{当检索器的服务对象从人转向 Agent，检索结果的消费方式与训练监督来源也随之变化\textsuperscript{[1]}。}
    \end{center}

    \vspace{0.5mm}
    \begin{center}
        \begin{tikzpicture}
            \node[font=\small\bfseries,text=navyBlue] at (-1.10,1.42) {Human Search};
            \node[font=\small\bfseries,text=navyBlue] at (5.20,1.42) {训练数据};

            \node[flow,minimum width=2.85cm,minimum height=0.72cm,font=\small] (hq) at (-4.30,0.62)
                {用户 query};
            \node[flowgreen,minimum width=2.85cm,minimum height=0.72cm,font=\small] (hr) at (-1.10,0.62)
                {排序结果};
            \node[flow,minimum width=2.85cm,minimum height=0.72cm,font=\small] (hf) at (2.10,0.62)
                {click / dwell};
            \draw[arrow] (hq) -- (hr);
            \draw[arrow] (hr) -- (hf);

            \node[card,minimum width=2.00cm,minimum height=0.72cm,font=\scriptsize] (hlog) at (5.35,0.62)
                {Human Log};

            \draw[->,line width=1.0pt,accentRed] (-1.10,0.22) -- (-1.10,-0.62);
            \node[font=\scriptsize\bfseries,text=accentRed,align=right,anchor=east]
                at (-1.34,-0.18) {检索结果不再是终点};
            \node[font=\scriptsize\bfseries,text=accentRed,align=left,anchor=west]
                at (-0.86,-0.18) {而是 Agent 后续推理的输入};

            \draw[->,line width=1.0pt,accentRed] (5.35,0.22) -- (5.35,-1.36);

            \node[font=\small\bfseries,text=navyBlue] at (-1.10,-0.88) {Agentic Search};
            \node[accentcard,minimum width=2.00cm,minimum height=0.72cm,font=\scriptsize] (atraj) at (5.35,-1.73)
                {Agent Trajectory};

            \node[flow,minimum width=2.30cm,minimum height=0.72cm,font=\scriptsize] (aq) at (-5.00,-1.73)
                {中间 query};
            \node[flowgreen,minimum width=2.30cm,minimum height=0.72cm,font=\scriptsize] (as) at (-2.45,-1.73)
                {snippet};
            \node[flow,minimum width=2.30cm,minimum height=0.72cm,font=\scriptsize] (ab) at (0.10,-1.73)
                {Browse};
            \node[flowgreen,minimum width=2.30cm,minimum height=0.72cm,font=\scriptsize] (ar) at (2.65,-1.73)
                {后续 reasoning};
            \draw[arrow] (aq) -- (as);
            \draw[arrow] (as) -- (ab);
            \draw[arrow] (ab) -- (ar);
        \end{tikzpicture}
    \end{center}

    \vspace{-1mm}
    {\centering\tiny\textcolor{mutedGray}{[1] Yuqi Zhou et al., \textit{Learning to Retrieve from Agent Trajectories}, SIGIR 2026.}\par}
    \vspace{-1mm}
    \begin{center}
        \compacttakeaway{如何利用 Agent Trajectory 训练面向 Agentic Search 的 retriever？}
    \end{center}
\end{frame}
